\documentclass[11pt]{article}
\usepackage[T1]{fontenc}
\usepackage{amsmath,amssymb,amsthm}
\usepackage[margin=1in]{geometry}
\usepackage{hyperref}
\hypersetup{hidelinks}
\newtheorem{theorem}{Theorem}[section]
\newtheorem{proposition}[theorem]{Proposition}
\newtheorem{lemma}[theorem]{Lemma}
\newtheorem{corollary}[theorem]{Corollary}
\theoremstyle{remark}
\newtheorem{remark}[theorem]{Remark}
\newcommand{\NN}{\mathbb N}
\newcommand{\ZZ}{\mathbb Z}
\newcommand{\RR}{\mathbb R}
\newcommand{\kk}{\Bbbk}
\newcommand{\one}{\mathbf 1}
\newcommand{\sdef}{\operatorname{sdef}}
\newcommand{\isdef}{\operatorname{isdef}}
\newcommand{\conv}{\operatorname{conv}}
\newcommand{\vol}{\operatorname{vol}}
\newcommand{\FRAC}{\operatorname{FRAC}}
\newcommand{\STAB}{\operatorname{STAB}}
\newcommand{\Hilb}{\operatorname{Hilb}}
\title{Symbolic defect and normalization of whiskered cover ideals}
\author{}
\date{September 8, 2026}
\begin{document}
\maketitle
\begin{abstract}
Let $J$ be the cover ideal of the full whiskering of a graph $G$ on $n$
vertices. We express its symbolic defect and integral symbolic defect as
differences between the Ehrhart functions of the fractional and ordinary
stable-set polytopes and the Hilbert function of the stable-set semigroup.
For every nonbipartite core graph the two defects are asymptotic, with common
leading term $(\vol(\FRAC(G))-\vol(\STAB(G)))m^n$ and positive coefficient.
Their difference is eventually polynomial of degree at most $n-1$, without
any normality hypothesis. This bound is sharp: when the complement of $G$
is a disjoint union of two odd cycles of lengths at least five, we determine
the difference in every degree by a single binomial coefficient, and give a
composition formula for any number of odd-cycle components. We also
give closed formulas for the full symbolic defect of whiskered complete
graphs.
\end{abstract}

\section{Introduction and statements}

For a homogeneous ideal $I$ in a polynomial ring, the symbolic defect
$\sdef(I,m)=\mu(I^{(m)}/I^m)$ measures the number of generators needed to
account for the difference between symbolic and ordinary powers. The
integral symbolic defect introduced by Oltsik \cite[Definition~4.1]{Oltsik}
is
\[
\isdef(I,m)=\mu\bigl(\overline{I^{(m)}}/\overline{I^m}\bigr).
\]
Here a bar denotes integral closure and $\mu$ the minimal number of module
generators. These two counts need not agree. Oltsik asks about their growth
comparison and the degrees and coefficients of their eventual
quasipolynomials \cite[Section~7]{Oltsik}. We obtain an exact description
for cover ideals of fully whiskered graphs, including graphs whose
stable-set semigroup rings are not normal.

We write $\NN=\ZZ_{\geq0}$. Throughout, $G$ is a finite simple graph with specified vertex set $[n]$,
where $n\geq1$. Isolated vertices are allowed. Its full whiskering $w(G)$
has vertices $x_1,\ldots,x_n,y_1,\ldots,y_n$, the core edges $x_ix_j$ for
$ij\in E(G)$, and the pendant edges $x_iy_i$. Over an arbitrary field
$\kk$, put
\[
R=\kk[x_1,\ldots,x_n,y_1,\ldots,y_n],\qquad
J=J(w(G))=\bigcap_{uv\in E(w(G))}(u,v).
\]
Write $s_m=\sdef(J,m)$ and $i_m=\isdef(J,m)$.

Let $\mathcal I(G)$ be the independent sets of $G$, including the empty
set. We use the polytopes and finite sets
\begin{align*}
P=\FRAC(G)&=\{b\in[0,1]^n:b_i+b_j\leq1\text{ for }ij\in E(G)\},\\
Q=\STAB(G)&=\conv\{\one_A:A\in\mathcal I(G)\},\\
B_m&=\{\one_{A_1}+\cdots+\one_{A_m}:A_j\in\mathcal I(G)\}.
\end{align*}
The upper coordinate bounds in $P$ include isolated vertices. Define
\[
E_P(m)=|mP\cap\ZZ^n|,\quad E_Q(m)=|mQ\cap\ZZ^n|,
\quad H_G(m)=|B_m|.
\]
All volumes below are Euclidean $n$-dimensional volumes, with unit cube
volume one. Thus the Ehrhart leading coefficient is $\vol(P)$, rather
than $n!\vol(P)$.

\begin{theorem}\label{thm:main}
For every graph $G$ as above and every $m\geq1$,
\begin{equation}\label{eq:dictionary}
s_m=E_P(m)-H_G(m),\qquad i_m=E_P(m)-E_Q(m).
\end{equation}
The difference $s_m-i_m$ is nonnegative, eventually polynomial, and
$O(m^{n-1})$. It vanishes for every $m$ if and only if $Q$ has the integer
decomposition property.

Both defects are eventually quasipolynomial with period dividing two.
If $G$ is bipartite, both vanish identically. If $G$ is nonbipartite, then
\begin{equation}\label{eq:asymptotic}
s_m=c_Gm^n+O(m^{n-1}),\qquad
i_m=c_Gm^n+O(m^{n-1}),\qquad
c_G=\vol(P)-\vol(Q)>0.
\end{equation}
In particular both parity branches have degree $n$ and leading coefficient
$c_G$, and $s_m/i_m\longrightarrow1$.
\end{theorem}

The next result gives an exact correction in a nonnormal family. For a
graph $G$, its complement is denoted by $\overline G$; this notation is
distinct from integral closure of ideals.

\begin{theorem}\label{thm:correction}
Suppose $\overline G=C_r\sqcup C_s$, where $r,s\geq5$ are odd, and put
$n=r+s$. Then
\begin{equation}\label{eq:correction}
s_m-i_m=
\begin{cases}
0,&m<n/2,\\[2pt]
\displaystyle\binom{m+n/2-1}{n-1},&m\geq n/2.
\end{cases}
\end{equation}
Consequently the exponent $n-1$ in Theorem~\ref{thm:main} is sharp. The
correction first occurs in degree $n/2$ and has leading coefficient
$1/(n-1)!$.
\end{theorem}

For complete cores, the normalization correction is zero and the full
defect has a direct closed form.

\begin{theorem}\label{thm:complete}
For $G=K_n$ one has $s_m=i_m$ for all $m$. For $q\geq0$,
\begin{align}
s_{2q}&=(q+1)^n+n\sum_{k=1}^q k^{n-1}-\binom{2q+n}{n},\label{eq:even}\\
s_{2q+1}&=(q+1)^n+n\sum_{k=1}^{q+1}k^{n-1}
-\binom{2q+1+n}{n}.\label{eq:odd}
\end{align}
In \eqref{eq:even}, $s_0$ is interpreted as zero. For $n\geq3$, the
leading coefficient as a function of $m$ is $2^{1-n}-1/n!>0$. For
$n=1,2$, all defects are zero.
\end{theorem}

\subsection*{Relation to previous work}
The ordinary generators of whiskered cover ideals, their independent-set
parametrization, and their analytic spread $n+1$ occur in
Drabkin--Guerrieri \cite[Proposition~5.7 and Theorem~5.8]{DGFreiman}.
The additional observation used here is that every symbolic power is
also generated in a single degree, and that ordinary products and integral
closure can be compared inside the same affine slice.

Eventual quasipolynomiality with period dividing two for graph cover
ideals is already known \cite[Corollary~4.3]{DG}. The conditional recursion
in \cite[Theorem~5.5]{DG} does not apply to all whiskered graphs:
\cite[Remark~5.3]{DG} exhibits the whiskered triangle as a failure of its
indecomposability hypothesis. Theorem~\ref{thm:complete} includes that
example and concerns whiskered complete graphs, whereas
\cite[Theorem~5.7]{DG} treats unwhiskered complete graphs.

The polyhedral facts about fractional stable-set polytopes used below
are established \cite{HHO}. Conditional comparisons with integral closure
appear in \cite[Proposition~4.5]{Oltsik} and, more recently,
\cite[Theorem~4.1]{BHL}; the latter paper also gives coefficient constancy
results under filtration and module hypotheses
\cite[Theorem~4.2]{BHL}. Our conclusions are the explicit volume difference,
unconditional ratio one for this whiskered class, and exact correction
\eqref{eq:correction}. They do not resolve the unrestricted growth
questions for monomial ideals.

Nonnormality in Theorem~\ref{thm:correction} is already implied by
\cite[Theorem~1]{MOS}. Preservation of multiplicity under semigroup
normalization is standard; see also \cite[Proposition~3.1]{Higashitani}.
Moreover, \cite[Proposition~4.2]{Higashitani} computes a translated free-face
hole monoid and its Hilbert-series correction for a different connected
edge-ring family built around two triangles. Our exact correction is a
specific stable-set and symbolic-defect calculation with that conceptual
precedent. We give elementary proofs, including the exclusion of holes
away from the relevant facet.

\section{The common degree slice}

Since $J$ is squarefree with minimal primes indexed by the edges,
\[
J^{(m)}=\bigcap_{uv\in E(w(G))}(u,v)^m.
\]
Thus $x^ay^b\in J^{(m)}$ if and only if
\begin{equation}\label{eq:covers}
a_i+b_i\geq m\quad(1\leq i\leq n),\qquad
a_i+a_j\geq m\quad(ij\in E(G)).
\end{equation}

\begin{lemma}\label{lem:generators}
The minimal monomial generators of $J^{(m)}$ are precisely
\[
u_b=\prod_{i=1}^n x_i^{m-b_i}y_i^{b_i},
\qquad b\in mP\cap\ZZ^n.
\]
In particular, every such generator has degree $nm$.
\end{lemma}
\begin{proof}
In a minimal generator, no $a_i$ can exceed $m$, because replacing it by
$m$ preserves every inequality in \eqref{eq:covers}. Once $a_i\leq m$,
a strict inequality $a_i+b_i>m$ permits decreasing the positive exponent
$b_i$. Hence minimality requires $a_i+b_i=m$ for every $i$.
Conversely, these equalities prevent decreasing any positive exponent:
its pendant inequality would fail. Substituting $a_i=m-b_i$ in the core
inequalities gives $b_i+b_j\leq m$, exactly the asserted lattice set.
\end{proof}

\begin{proposition}\label{prop:dictionary}
The identities \eqref{eq:dictionary} hold for every $m\geq1$.
\end{proposition}
\begin{proof}
At $m=1$, the lattice vectors in Lemma~\ref{lem:generators} are precisely
the independent-set indicators. A product of $m$ such generators is
$u_b$ with $b\in B_m$, and every element of $B_m$ occurs. These products
all have degree $nm$, so distinct products are identified only when their
exponent vectors are equal.

Let $\mathfrak m$ be the homogeneous maximal ideal of $R$. If a monomial
ideal $A$ is generated in one degree $d$ and $D\subseteq A$ is a monomial
ideal, the minimal generators of $A/D$ are exactly the degree-$d$
generators of $A$ absent from $D$. Their nonzero classes form a basis
modulo $\mathfrak m(A/D)$: the latter has no degree-$d$ part, and the
monomials present in the quotient are linearly independent. Applying this
to $A=J^{(m)}$, $D=J^m$ gives $s_m=E_P(m)-H_G(m)$.

For a monomial ideal, a monomial belongs to its integral closure precisely
when its exponent lies in its Newton polyhedron. This criterion can also
be expressed by requiring a positive power of the monomial to belong to
the same power of the ideal; see \cite[Proposition~2.13]{Oltsik}.
The Newton polyhedron of $J^m$ is the convex hull of its generating
exponent vectors plus $\RR_{\geq0}^{2n}$. Its intersection with the
total-degree hyperplane $nm$ is exactly the affine image of $mQ$ under
\begin{equation}\label{eq:affine}
b\longmapsto(m\one-b,b).
\end{equation}
Indeed the convex hull of all sums of $m$ independent-set indicators is
$mQ$. All these exponent vectors already have total degree $nm$, and a
nonzero nonnegative addition would increase that degree. It follows that
\[
u_b\in\overline{J^m}\quad\Longleftrightarrow\quad b\in mQ.
\]
Each coordinate-prime power $(u,v)^m$ is integrally closed, as is their
intersection. Thus $J^{(m)}$ is integrally closed and
$\overline{J^m}\subseteq J^{(m)}$. Applying the same degree-$nm$ quotient
argument gives $i_m=E_P(m)-E_Q(m)$. This argument does not assume that the
entire ideal $\overline{J^m}$ is generated in degree $nm$.
\end{proof}

\section{Normalization and asymptotic comparison}

We include the semigroup estimate to make clear that no normality
assumption on the stable-set ring is needed.

\begin{lemma}\label{lem:conductor}
Let $S\subseteq\ZZ^{n+1}$ be generated by $(\one_A,1)$ for
$A\in\mathcal I(G)$, and let
$\overline S=\RR_{\geq0}S\cap\ZZ^{n+1}$. Then
\[
0\leq E_Q(m)-H_G(m)=O(m^{n-1}).
\]
The difference is eventually a polynomial.
\end{lemma}
\begin{proof}
The group generated by $S$ is all of $\ZZ^{n+1}$: it contains
$(0,1)$ and $(e_i,1)$, and their differences give $(e_i,0)$.
The degree-$m$ elements of $\overline S$ are exactly
$\{(b,m):b\in mQ\cap\ZZ^n\}$.

There exists a conductor element $c\in S$ with $c+\overline S\subseteq S$.
Here is a direct construction. If $a_1,\ldots,a_t$ generate $S$, subtract
the floors of the coefficients in a nonnegative real expression
$u=\sum_j\lambda_ja_j$ for $u\in\overline S$. The remainder is a lattice
point in the bounded set $\{\sum_j\theta_ja_j:0\leq\theta_j<1\}$.
Its finitely many possible values $r_1,\ldots,r_v$ belong to
$\overline S$, and $\overline S\subseteq\bigcup_j(r_j+S)$.
Fullness of the generated group allows us to write
$r_j=p_j-q_j$ with $p_j,q_j\in S$. Taking $c=\sum_jq_j$ ensures
$c+r_j\in S$ for every $j$, as required.

Write $c=(c',d)$. Translation by $c$ injects
$\overline S_{m-d}$ into $S_m$ for $m\geq d$, and hence
\[
0\leq E_Q(m)-H_G(m)\leq E_Q(m)-E_Q(m-d)=O(m^{n-1}).
\]
The last assertion uses Ehrhart's theorem: $Q$ is an integral,
full-dimensional polytope, containing $0,e_1,\ldots,e_n$, so $E_Q$ is
a polynomial of degree $n$. Finally $H_G$ is the Hilbert function of the
standard graded affine semigroup algebra $\kk[S]$, and is therefore
eventually polynomial by the Hilbert--Serre theorem. Their difference is
eventually polynomial as well.
\end{proof}

\begin{lemma}\label{lem:half}
Every vertex of $P$ has coordinates in $\{0,1/2,1\}$.
\end{lemma}
\begin{proof}
Let $b$ be a vertex. Consider the graph on coordinates with
$0<b_i<1$, retaining the edges on which $b_i+b_j=1$. A tight edge cannot
join a fractional coordinate to a coordinate equal to zero or one.
If a connected component of this tight graph is bipartite, add
$\varepsilon$ on one side and subtract it on the other. For sufficiently
small positive and negative $\varepsilon$, every nontight inequality and
every coordinate bound remains satisfied. This contradicts extremality.
Every component therefore has an odd cycle. The tight equations around
that cycle force all its coordinates to be $1/2$, and the same equations
propagate this value throughout the component.
\end{proof}

\begin{proof}[Proof of Theorem~\ref{thm:main}]
Proposition~\ref{prop:dictionary} and Lemma~\ref{lem:conductor} give the
identities and correction estimate. Since $Q\subseteq[0,1]^n$ satisfies
the edge inequalities, its lattice points are exactly the independent-set
indicators. Thus $E_Q(m)=H_G(m)$ for every $m$ is precisely the integer
decomposition property of $Q$.

By Lemma~\ref{lem:half} and Ehrhart's theorem, $E_P$ is a
quasipolynomial of period dividing two, with both branches having leading
term $\vol(P)m^n$. Also
\[
E_Q(m)=\vol(Q)m^n+O(m^{n-1}),\qquad
H_G(m)=\vol(Q)m^n+O(m^{n-1}).
\]
The asserted eventual quasipolynomiality and asymptotic expansions follow.

Suppose first that $G$ is bipartite with parts $U,V$. For
$b\in mP\cap\ZZ^n$, assign vertex $i\in U$ the first $b_i$ colors from
$[m]$, and vertex $j\in V$ the last $b_j$ colors. Adjacent vertices
receive disjoint color sets because $b_i+b_j\leq m$. Each color class is
therefore independent, giving a representation of $b$ in $B_m$.
Isolated vertices can be placed in either part. Hence $B_m=mP\cap\ZZ^n$
and both defects vanish.

If $G$ is not bipartite, choose an odd cycle with $2r+1$ vertices. On its
vertex set the coordinate sum of any point of $Q$ is at most $r$. The
point $(1/2,\ldots,1/2)$ belongs to $P$ and violates that inequality, so
$Q\subsetneq P$. Both polytopes are full-dimensional. Their strict
containment gives $\vol(Q)<\vol(P)$: moving a point of $P\setminus Q$
slightly toward an interior point of $P$ produces a small open ball in
$P$ disjoint from the closed set $Q$. Thus $c_G>0$. It follows that both
parity branches have degree exactly $n$, and that $i_m>0$ eventually.
Dividing the two asymptotic expansions proves the ratio limit.
\end{proof}

\section{An exact correction from two odd cycles}

Let $C$ be an odd cycle and $b$ a nonnegative integral demand vector on
its vertices. Put $T=\sum_i b_i$ and define
\begin{align*}
\nu^*(b)&=\max\left\{\sum_{e\in E(C)}\lambda_e:
\lambda_e\geq0,\ \sum_{e\ni i}\lambda_e\leq b_i\right\},\\
\nu(b)&=\max\left\{\sum_{e\in E(C)}\lambda_e:
\lambda_e\in\NN,\ \sum_{e\ni i}\lambda_e\leq b_i\right\}.
\end{align*}
Repeated use of an edge is allowed in this capacitated matching problem.

\begin{lemma}\label{lem:cycle}
One has $\nu(b)=\lfloor\nu^*(b)\rfloor$. A noninteger value of
$\nu^*(b)$ occurs exactly when the full saturation equations
$\lambda_{i-1}+\lambda_i=b_i$ have a nonnegative, nonintegral solution.
In that case the solution is unique, every edge weight is a positive
halfinteger, $T$ is odd, and $\nu^*(b)=T/2$.
\end{lemma}
\begin{proof}
Choose an optimal vertex of the bounded fractional matching polytope.
If an edge coordinate is zero, its positive support is a disjoint union
of paths. The vertex-edge incidence matrix of a bipartite graph is
totally unimodular: changing signs on the rows of one bipartition gives
an oriented incidence matrix, whose square subdeterminants are
$0,1,-1$. Thus the resulting active system with integral right-hand
side has an integral vertex.

If every edge coordinate is positive, extremality requires every vertex
inequality to be tight; otherwise there are fewer active constraints
than edge variables. The incidence matrix of an odd cycle is
nonsingular, as its homogeneous equations alternate signs and the odd
closing condition forces zero. Solving successively shows that its
saturated solution is either integral or has every coordinate in
$\ZZ+1/2$. In the latter case all coordinates are at least $1/2$,
and $T$ is odd. Round consecutive edge coordinates alternately down,
up, down, up, ending down. All vertex sums remain unchanged except
the one between the first and last edges, where the sum decreases by
one. The resulting integral matching has weight $(T-1)/2$.

This proves the floor identity. Conversely, any feasible nonintegral
saturated solution attains the universal upper bound $T/2$, and its
uniform halfintegrality on an odd number of edges makes $T$ odd.
\end{proof}

\begin{proof}[Proof of Theorem~\ref{thm:correction}]
Set $F=\overline G=C_r\sqcup C_s$. Because both cycles have length at
least five, the independent sets of $G$ are exactly the empty set,
singletons, and edges of $F$. For $b\in\NN^n$, put $T=\sum_i b_i$,
and let $\nu_j^*,\nu_j$ be the two matching numbers on component $j$.
Then
\begin{align}
b\in mQ&\quad\Longleftrightarrow\quad
T-\nu_1^*(b)-\nu_2^*(b)\leq m,\label{eq:fractional-test}\\
b\in B_m&\quad\Longleftrightarrow\quad
T-\nu_1(b)-\nu_2(b)\leq m.\label{eq:integer-test}
\end{align}
Indeed a choice of edge weights leaves nonnegative residual demands
filled by singleton weights. The total coefficient mass is
$T-\sum_e\lambda_e$; empty-set weights fill any remaining mass up to
$m$. This argument works over the nonnegative reals for
\eqref{eq:fractional-test} and over the nonnegative integers for
\eqref{eq:integer-test}.

If both fractional optima are integral, the two tests coincide. If just
one is a halfinteger, rounding the fractional threshold up to the
integer $m$ again gives the ordinary threshold. Thus a hole
$b\in(mQ\cap\ZZ^n)\setminus B_m$ can occur only when both fractional
optima are halfintegers. By Lemma~\ref{lem:cycle}, each component is then
fully saturated by its unique positive halfintegral edge vector, and
its demand total is odd. The fractional threshold is $T/2$, an integer,
whereas the ordinary threshold is $T/2+1$. Consequently a hole occurs
exactly when $T=2m$ and both saturation solutions are positive
halfintegral. This also excludes holes away from the face
$\sum_i b_i=2m$.

Write these edge weights as $\lambda_e=k_e+1/2$, with $k_e\in\NN$.
The block incidence matrix of the two odd cycles is nonsingular, so
edge weights determine demands injectively. Conversely every such
edge vector gives integral nonnegative demands, odd component totals,
and the hole condition at degree $m=\sum_e\lambda_e$. We have obtained
a bijection between degree-$m$ holes and
\[
\left\{(k_e)_{e\in E(F)}\in\NN^n:
\sum_e k_e=m-\frac n2\right\}.
\]
Stars and bars gives \eqref{eq:correction}, by
Theorem~\ref{thm:main}. Its first degree and leading coefficient follow
immediately.
\end{proof}

\begin{remark}
Equivalently, the Hilbert-series correction for this family is
\[
\sum_{m\geq0}(s_m-i_m)t^m=\frac{t^{n/2}}{(1-t)^n}.
\]
The correction consists of a translate of the free monoid on the $n$
edge vectors of $C_r\sqcup C_s$. For $r=s=5$, its first values are
$0,0,0,0,1,10,55$ in degrees $1$ through $7$.
\end{remark}

The same threshold argument gives a composition formula with any number
of odd-cycle components. The factors below are the Hilbert series of
single-component stable-set rings, so the formula uses ordinary Hilbert
data to recover the entire normalization correction of their graph join.

\begin{proposition}\label{prop:composition}
Suppose $\overline G=\bigsqcup_{a=1}^t C_{r_a}$, with every $r_a\geq5$
odd. Let
\begin{align*}
\mathcal H_a(z)&=\sum_{m\geq0}H_{\overline{C_{r_a}}}(m)z^m,
&V_a(z)&=\frac{z^{(r_a+1)/2}}{(1-z)^{r_a}},\\
U_a(z)&=(1-z)\mathcal H_a(z)-V_a(z).
\end{align*}
Then
\begin{equation}\label{eq:composition}
\sum_{m\geq0}(s_m-i_m)z^m
=\sum_{\substack{F\subseteq[t]\\ |F|\geq2}}
\left(\sum_{j=1}^{\lfloor|F|/2\rfloor}z^{-j}\right)
\prod_{a\in F}V_a(z)\prod_{a\notin F}U_a(z).
\end{equation}
The right-hand side is a formal power series with nonnegative exponents.
For $t=2$ it equals
\[
\frac{z^{(r_1+r_2)/2}}{(1-z)^{r_1+r_2}}.
\]
\end{proposition}
\begin{proof}
For a demand vector on component $a$, let
$c_a=T_a-\nu_a$ be its minimum ordinary coefficient mass, and mark the
component if $\nu_a^*$ is a halfinteger. By Lemma~\ref{lem:cycle}, the
minimum fractional mass is $c_a$ on an unmarked component and
$c_a-1/2$ on a marked one. In particular the single-component polytope
has the integer decomposition property: for an integer degree, its
fractional and ordinary feasibility thresholds agree.

The series counting all component demands by their minimum ordinary
mass is $(1-z)\mathcal H_a(z)$, since the Hilbert function counts
demands whose minimum mass is at most the given degree. A marked
component has unique edge weights $k_e+1/2$ and
\[
c_a=\sum_e k_e+\frac{r_a+1}{2}.
\]
Thus $V_a$ counts the marked demands by $c_a$, and $U_a$ counts the
unmarked demands. All coefficients are nonnegative and each
coefficient is finite, because total demand is at most twice the
minimum mass.

For a demand on all components, put $C=\sum_a c_a$ and let $F$ be the
set of marked components. Independent sets of $G$ lie within a single
component of its complement, so the ordinary minimum mass is $C$ and
the fractional minimum mass is $C-|F|/2$. Its missing degrees are
exactly $C-j$, for $1\leq j\leq\lfloor|F|/2\rfloor$. Multiplying the
component demand series and recording those degrees gives
\eqref{eq:composition}. Each factor $V_a$ starts in degree at least
three, so the displayed negative powers cannot create negative
exponents. The specialization $t=2$ follows directly.
\end{proof}

\section{Complete cores}

\begin{proof}[Proof of Theorem~\ref{thm:complete}]
For $K_n$, independent sets are empty or singletons. Hence
\[
B_m=\{b\in\NN^n:\textstyle\sum_i b_i\leq m\},\qquad
H_G(m)=E_Q(m)=\binom{m+n}{n}.
\]
It remains to count integral $b$ with $0\leq b_i\leq m$ and
$b_i+b_j\leq m$ for distinct $i,j$. Put $q=\lfloor m/2\rfloor$.
If every coordinate is at most $q$, there are $(q+1)^n$ choices.
Otherwise exactly one coordinate exceeds $q$. If that coordinate is
$a$, every other coordinate ranges independently from $0$ to $m-a$;
their mutual pair constraints are automatic. The additional count is
\[
n\sum_{a=q+1}^m(m-a+1)^{n-1}.
\]
Separating even and odd $m$ and subtracting $\binom{m+n}{n}$ proves
\eqref{eq:even}--\eqref{eq:odd}. The same partition works for $n=1$,
where the coordinate bound is retained even though there are no edges.
For $n=1,2$ the graph is bipartite, so the defect vanishes.
Finally, for $n\geq3$ the power-sum estimate
$\sum_{k=1}^q k^{n-1}=q^n/n+O(q^{n-1})$ gives leading coefficient
$2^{1-n}-1/n!$. It is positive because $n!>2^{n-1}$ for $n\geq3$.
\end{proof}

In particular the whiskered triangle has $s_2=1$ and $s_3=3$. The
calculation includes the obstruction in \cite[Remark~5.3]{DG} without
using the indecomposability recursion.

\subsection*{Verification and scope}
The proofs above are independent of computation. Ancillary scripts
check direct minimality of symbolic exponent vectors on all core graphs
with at most three vertices and powers at most three; ordinary and
symbolic counting on all $1099$ labeled graphs with at most five
vertices in powers one through three; and the complete-core formulas
for $n,m\leq6$. A separate check compares fractional matching vertex
enumeration with exact integer matching dynamic programming on $11628$
demand vectors on $C_5$, and verifies the two-component correction
through degree seven. All checks pass. The fractional vertex calculation
uses numerical linear algebra with a checked halfinteger tolerance;
Lemma~\ref{lem:cycle} supplies the exact argument.

The results concern the fully whiskered cover-ideal class. They neither
classify normal stable-set polytopes nor determine the minimal
quasiperiod for arbitrary cores. The complete-core formula and the
nonnormal correction family are parts of the same comparison theorem.

\begin{thebibliography}{9}

\bibitem{BHL}
A.~Banerjee, T.~H.~H\`a, and V.~B.~Lama,
\emph{Defect Functions Between Filtrations of Ideals},
preprint (2025),
\href{https://arxiv.org/abs/2512.14573}{arXiv:2512.14573}.

\bibitem{DG}
B.~Drabkin and L.~Guerrieri,
\emph{Asymptotic invariants of ideals with Noetherian symbolic Rees
algebra and applications to cover ideals},
J. Pure Appl. Algebra \textbf{224} (2020), 300--319.
\href{https://doi.org/10.1016/j.jpaa.2019.05.008}{doi:10.1016/j.jpaa.2019.05.008}.

\bibitem{DGFreiman}
B.~Drabkin and L.~Guerrieri,
\emph{On quasi-equigenerated and Freiman cover ideals of graphs},
Comm. Algebra \textbf{48} (2020), no.~10, 4413--4435.
\href{https://doi.org/10.1080/00927872.2020.1762890}{doi:10.1080/00927872.2020.1762890}.

\bibitem{HHO}
G.~Hamano, T.~Hibi, and H.~Ohsugi,
\emph{Ehrhart series of fractional stable set polytopes of finite graphs},
Ann. Comb. \textbf{22} (2018), 563--573.
\href{https://doi.org/10.1007/s00026-018-0392-2}{doi:10.1007/s00026-018-0392-2}.

\bibitem{Higashitani}
A.~Higashitani,
\emph{Difference of Hilbert series of homogeneous monoid algebras and
their normalizations},
Semigroup Forum \textbf{108} (2024), 101--114.
\href{https://doi.org/10.1007/s00233-024-10414-0}{doi:10.1007/s00233-024-10414-0}.

\bibitem{MOS}
K.~Matsuda, H.~Ohsugi, and K.~Shibata,
\emph{Toric Rings and Ideals of Stable Set Polytopes},
Mathematics \textbf{7} (2019), no.~7, article 613.
\href{https://doi.org/10.3390/math7070613}{doi:10.3390/math7070613}.

\bibitem{Oltsik}
B.~R.~Oltsik,
\emph{Symbolic Defect of Monomial Ideals},
Comm. Algebra \textbf{52} (2024), no.~9, 3996--4012;
\href{https://arxiv.org/abs/2310.12280}{arXiv:2310.12280}.

\end{thebibliography}
\end{document}
